\uuid{I057}
\titre{Primitive}
\auteur{bodin ; croizat ; sacre}
\organisation{exo7}
\licence{CC BY-NC-SA 4.0}
\datecreate{2022-03-08}
\niveau{L1}
\module{Analyse}
\chapitre{Calcul d'intégrales}
\sousChapitre{Primitives diverses}
\difficulte{2}
\theme{Primitive}
\exo7qcmid{hilisit.equations-differentielles.006}

\contenu{
\begin{question}
\qclassification{127.11}
Quelles sont les affirmations vraies ?
\begin{answers}
  \bad{$x\mapsto \ln(x)$ est une primitive de $x\mapsto 1/x$ sur $\Rr$.}
  \bad{$x\mapsto \ln(x)$ est une primitive de $x\mapsto 1/x$ sur $]-\infty,0[$.}
  \good{$x\mapsto \ln(x)$ est une primitive de $x\mapsto 1/x$ sur $]0,+\infty[$.}
  \good{$x\mapsto \ln(-x)$ est une primitive de $x\mapsto 1/x$ sur $]-\infty,0[$.}
\end{answers}
\begin{explanations}
La fonction $\ln$ n'est définie et dérivable que sur $]0,+\infty[$. Pour tout $x$ de $]0,+\infty[$, $(\ln(x))'=1/x$ ; pour tout $x$ de $]-\infty,0[$, la fonction $x \mapsto \ln(-x)$ est bien définie et dérivable, et on a $(\ln(-x))'=-1/-x=1/x$.
\end{explanations}
\end{question}
}
